% JCAP submission — initial-submission preprint format.
% NOTE: JCAP's official class is jcappub.sty (SISSA/IOP template); it is not
% installed on this machine, so this manuscript is prepared in the article
% class with JCAP-compatible structure (title/author/abstract/keywords,
% numbered sections, numeric citations). JCAP accepts an initial PDF
% submission in any single-column format; reformat to jcappub at the
% revision stage if requested.
\documentclass[11pt,a4paper]{article}
\usepackage[T1]{fontenc}
\usepackage[utf8]{inputenc}
\usepackage{lmodern}
\usepackage{amsmath,amssymb}
\usepackage{booktabs}
\usepackage{geometry}
\geometry{margin=2.7cm}
\usepackage{microtype}
\usepackage[hidelinks]{hyperref}
\newcommand{\Mbar}{M_{\rm bar}}
\newcommand{\rhoDE}{\rho_{\rm DE}}
\newcommand{\azero}{a_0}
\newcommand{\msms}{{\rm m\,s^{-2}}}

\title{\bfseries A parameter-free, pre-registered test that the galaxy
acceleration scale tracks the dark-energy density:\\
$a_0(z)/a_0(0)=\sqrt{\rho_{\rm DE}(z)/\rho_{\rm DE,0}}$}

\author{Carl P. Zimmerman\\[2pt]
\normalsize Briar Creek Tech (independent researcher)\\
\normalsize \href{mailto:carl@briarcreektech.com}{carl@briarcreektech.com}
\;\textperiodcentered\;
ORCID \href{https://orcid.org/0009-0008-3508-7982}{0009-0008-3508-7982}}
\date{July 2026}

\begin{document}
\maketitle

\begin{abstract}
\noindent
The galaxy acceleration scale $\azero\approx1.2\times10^{-10}\,\msms$ that sets
the onset of the mass-discrepancy/MOND regime coincides numerically with
$c\sqrt{\rho_\Lambda}$ --- a fact noted by Limbach, Psaltis \& \"Ozel (2008).
In the de Sitter--Unruh modified-inertia framework this scale is sourced by the
cosmological horizon, $\azero=cH_\Lambda/Z$; promoting it adiabatically as the
horizon evolves ties the coupling, under the adiabatic-horizon ansatz (itself
posited, not derived), to the dark-energy density and yields a parameter-free
redshift law, $\azero(z)/\azero(0)=\sqrt{\rhoDE(z)/\rho_{\rm DE,0}}$. We study
the ratio $R\equiv \azero(z)/\azero(0)$, which is independent of the
framework's internal normalization $Z$ and of the present-day value of
$\azero$ (both cancel), so that it depends only on the measured dark-energy
equation of state $w(z)$. This makes the framework falsifiable: $R$ tracks
$\sqrt{\rhoDE(z)}$; under the DESI-preferred ($w_0>-1$, $w_a<0$) histories it
declines at high redshift, it reduces to $R=1$ exactly if $w=-1$ (constant-$\azero$
MOND), and a measured rise registers as a strained verdict in the frozen
protocol. We confront the prediction non-circularly across two scales --- a
cosmology-side gate driven by DESI DR2 $w_0w_a$ constraints, which fixes the
required sign and amplitude, and an independent galaxy-side measurement of
$\azero(z)$ from the deep-MOND baryonic Tully--Fisher (BTFR) zero point. We
state the limitations plainly: the cosmology gate inherits $w(z)$ and is a
consistency condition rather than independent evidence; the reference ratio at
$z=3$ is a CPL extrapolation beyond the supernova range; the predicted
$\azero(z)$ is non-monotonic (a few-percent rise at $z\lesssim0.5$ before
declining at high $z$), so the two scales probe opposite branches of the same
curve; and the galaxy-side test is at present $\Mbar$-dominated and
underpowered (neither passed nor failed). The central methodological finding
redirects the observational strategy: in the BTFR estimator $\azero$ is
$1{:}1$ degenerate with $\Mbar$, the dominant $\Mbar$ systematic
(IMF, $\alpha_{\rm CO}$, dust) is \emph{common-mode} and does not average down
with sample size --- with a $\gtrsim0.08$-dex correlated floor the BTFR test
never reaches $3\sigma$ at any rotator count. The decisive lever is
common-mode $\Mbar$ calibration to $\lesssim0.05$ dex (with $\sim$20--40 clean
deep-MOND rotators), and the degeneracy-breaking observable is the
rotation-curve/radial-acceleration-relation transition, not the BTFR zero
point. At present DESI DR2 gives $R(z{=}3)=0.775\,[0.68,0.88]$, a $2.0\sigma$
hint below our committed $3\sigma$ bar --- undecided. We therefore
pre-register --- with a hashed, timestamped estimator, decision thresholds,
and data hierarchy fixed before the data exist --- the test against the
forthcoming Rubin/LSST calibrated supernova cosmology sample (expected
$\gtrsim2027$), which forecasts project to settle the gate at the $3$--$6\sigma$
level by $\sim$2028--2030. The value of $\azero$, the normalization $Z$, and
the sign of the effect are posited, not derived; the kernel
$\nu=\sqrt{1+1/y}$ is Milgrom's (1999) and the covariant MOND-scale
realization is AeST (Skordis \& Zlosnik 2021). The contribution is a bounded,
falsifiable, pre-registered program, not a detection.
\end{abstract}

\medskip
\noindent\textbf{Keywords:} modified gravity, dark energy equation of state,
supernova type Ia --- standard candles, redshift surveys, rotation curves of
galaxies.

\tableofcontents

\section{Introduction}
\label{sec:intro}

\subsection{The coincidence, and what promoting it to a law buys}

A spiral galaxy's mass discrepancy switches on below a universal acceleration
$\azero\approx1.2\times10^{-10}\,\msms$ (Milgrom 1983 \cite{Milgrom1983}). To
order unity this number equals the acceleration scale set by the dark energy
accelerating the whole universe, $c\sqrt{\rho_\Lambda}$ --- a coincidence that
is \emph{not} original to this work: Limbach, Psaltis \& \"Ozel (2008)
\cite{LPO2008} identified it and tested candidate $\azero$--$H_0$ and
$\azero$--$\rhoDE$ couplings against Tully--Fisher data to $z\simeq1.2$. If
the coincidence is physical --- if $\azero$ is genuinely sourced by the dark
sector --- then as the dark-energy density changes through cosmic history,
$\azero$ must change with it, on a track fixed by the cosmology with nothing
to tune. This paper states that track precisely, shows what current data do to
it, and pre-registers the decision procedure by which forthcoming data will
judge it.

\subsection{The framework and the horizon scale}
\label{sec:framework}

Modified inertia (MI) proposes that the inertial response of a body weakens
below a critical acceleration $\azero$, rather than the gravitational field
being modified. In the de Sitter--Unruh realization, an accelerated observer
in a universe with a cosmological horizon sees a floor in the vacuum's
response set by the horizon temperature, and the transition scale is tied to
the horizon:
\begin{equation}
\azero=\frac{cH_\Lambda}{Z}=\frac{c^2\sqrt{\Lambda/3}}{Z},
\qquad Z=\sqrt{\tfrac{32\pi}{3}}=5.78881,
\label{eq:a0}
\end{equation}
with $H_\Lambda\equiv c\sqrt{\Lambda/3}$ the de Sitter (pure dark-energy)
Hubble rate. The Planck 2018 value $\Lambda=1.089\times10^{-52}\,{\rm m^{-2}}$
gives the \textbf{canonical footing}
$\azero^{\rm canon}=c^2\sqrt{\Lambda/32\pi}=9.355\times10^{-11}\,\msms$ (pure
$\rhoDE$). A second \textbf{ALT footing} replaces the pure-$\Lambda$ density
by the total present-day critical density,
$\azero^{\rm ALT}=cH_0/Z=1.1305\times10^{-10}\,\msms$. The two differ by
$20.9\%$; both are carried on every dimensional number below (they cancel
identically in every ratio).

\paragraph{Scope, stated once and binding throughout.} Nothing here derives
the \emph{value} of $\azero$, the constant $Z$, or the \emph{sign} of the
inertial correction: those are posited inputs. The low-acceleration
interpolation used to read $\azero$ from rotation curves,
$\nu(y)=\sqrt{1+1/y}$ with $y=g_{\rm bar}/\azero$, is the identical functional
form to Milgrom's 1999 interpolation (\cite{Milgrom1999}, Eq.~9, in the
original MOND context of \cite{Milgrom1983}); Milgrom's coefficient was
$2cH_\Lambda$ rather than $cH_\Lambda/Z$. The covariant modified-\emph{gravity}
realization of a MOND scale is AeST \cite{AeST}; this work adopts both and
claims neither. The distinctive content claimed here is only the
$cH_\Lambda/Z$ \emph{coefficient}, the modified-inertia reading, and the
falsifiable $\azero(z)\propto\sqrt{\rhoDE(z)}$ law under test.

\subsection{Related work by the author}
\label{sec:related}

This manuscript consolidates, into one peer-review-facing article, three of
the author's Zenodo preprint/pre-registration deposits released for
reproducibility and timestamping: the parameter-free $\azero(z)$ law and its
DESI confrontation \cite{Zdesi}, the non-circular cross-scale test
\cite{Zcross}, and the Rubin/LSST pre-registration \cite{Zprereg}. Two
companion papers by the author treat logically distinct results and are cited,
not consolidated: a field-theory paper (covariant modified-inertia
construction and an obstruction theorem; Zenodo record \cite{Zmift}, PRD
submission in preparation) and a galaxy-dynamics data paper (the SPARC
$\azero$-line slope measurement and its $\Lambda$-inversion; Zenodo record
\cite{Za0line}, MNRAS submission in preparation). On the footing fork the
division of labor is explicit: \emph{the value-normalization fork (canonical
$cH_\Lambda/Z$ vs $cH_0/Z$) is examined empirically in the companion SPARC
measurement \cite{Za0line}; here only the evolution-law footing is under
test.} Nothing in the present test depends on the covariant field-theory
completion (Section~\ref{sec:limitations}).

\section{The prediction}
\label{sec:prediction}

\subsection{The parameter-free replacement of the free-$w$ fit}

The framework keeps General Relativity for the cosmological background and
therefore has no native supernova formula: light-curve-to-distance
standardization is unchanged, and the Friedmann equation is the standard one.
Its single distinctive move is that the dark-energy term is not a free fit but
the galaxy acceleration scale,
\begin{equation}
H^2(z)=\frac{8\pi G}{3}\,\rho_m(z)+\frac{Z^2\azero^2(z)}{c^2},
\qquad \frac{Z^2\azero^2}{c^2}=H_\Lambda^2=\frac{\Lambda c^2}{3}\;\;(z=0),
\end{equation}
so that, inverting $Z^2\azero^2/c^2=\tfrac{8\pi G}{3}\rhoDE$ with
$Z^2=32\pi/3$, the dark-energy mass density is
$\rhoDE(z)=4\azero^2(z)/(Gc^2)$ and
\begin{equation}
\boxed{\;\azero(z)=\tfrac{c}{2}\sqrt{G\,\rhoDE(z)}
=c^2\sqrt{\Lambda_{\rm eff}(z)/32\pi}
\quad\Longrightarrow\quad
R(z)\equiv\frac{\azero(z)}{\azero(0)}=\sqrt{\frac{\rhoDE(z)}{\rho_{\rm DE,0}}}\;}
\label{eq:law}
\end{equation}
The SI form $\tfrac{c}{2}\sqrt{G\rhoDE}$ uses the dark-energy \emph{mass}
density (kg\,m$^{-3}$); the equivalent geometric form uses
$\Lambda_{\rm eff}\equiv8\pi G\rhoDE/c^2$ (m$^{-2}$). Both give
$\azero^{\rm canon}=9.355\times10^{-11}\,\msms$ at $z=0$ on Planck
$\rho_{\rm DE,0}$; the two conventions must not be mixed. The
redshift-tracking ratio $R(z)$ is convention-independent.

On the supernova side this adds \textbf{no free dark-energy parameter} ---
against $\Lambda$CDM's $\Omega_\Lambda$ (and, in the evolving-DE extension,
$w_0,w_a$). The framework does not, however, produce a parameter-free
dark-energy history from nothing: it \emph{consumes} whatever $\rhoDE(z)$ the
cosmic data prefer (itself the $w_0,w_a$ fit to distances) and requires that
the \emph{same} function reproduce the galaxy-side $\azero(z)$. The content is
therefore a cross-\emph{dataset} constraint, not a first-principles prediction
of $w(z)$. At $z=0$ the framework is identical to $\Lambda$CDM by
construction. It differs from $\Lambda$CDM only if $\azero(z)$ evolves.

\subsection{$R$ is $Z$-independent, and the two distinct footing forks}
\label{sec:forks}

Because $R$ is a ratio of $\azero$ at two redshifts, the normalization $Z$ and
the absolute value of $\azero$ cancel identically: both \emph{value} footings
(canonical $9.36\times10^{-11}$ and ALT $1.13\times10^{-10}\,\msms$) collapse
onto the same $R(w_0,w_a)$. $R$ depends on the dark-energy expansion history
alone. One clarification is owed, because the framework's internal audit
distinguishes two separate footing forks. The fork that cancels is the
\emph{value} fork (which absolute $\azero$ today). A second, logically
independent fork concerns the \emph{evolution law} itself: the
adiabatic-horizon $\azero\propto\sqrt{\rhoDE(z)}$ (which declines at high $z$
under evolving DE) versus a rival $\azero\propto cH(z)=cH_0E(z)$ (which
\emph{rises} with $z$). \textbf{This work commits to the $\sqrt{\rhoDE}$
footing as the prediction under test}, in advance: if the data instead showed
a rise, that registers as Verdict~C (framework strained) in the frozen
protocol of Section~\ref{sec:gate} --- it is not re-interpreted post hoc as
the rival footing succeeding.

\subsection{What the framework predicts, and what it inherits}
\label{sec:tracks}

For the CPL parameterization $w(a)=w_0+w_a(1-a)$,
\begin{equation}
\frac{\rhoDE(z)}{\rho_{\rm DE,0}}=(1+z)^{3(1+w_0+w_a)}
\exp\!\left(\frac{-3w_a z}{1+z}\right).
\label{eq:cpl}
\end{equation}
Two reference cosmic tracks bracket the question, with the rival
evolution model shown for contrast:

\begin{center}
\begin{tabular}{lccc}
\toprule
track & $(w_0,\,w_a)$ & $\azero(3)/\azero(0)$ & $\azero(3.25)/\azero(0)$\\
\midrule
flat-$\Lambda$ ($w=-1$) & $(-1.00,\;0.00)$ & $1.000$ (flat) & $1.000$\\
DESI-DR2 evolving & $(-0.75,\;-0.90)$ & $0.712$ (declining) & $0.685$\\
\emph{(rival $cH(z)$ evolution model)} & $\azero\propto cH_0E(z)$ & $4.57$ (rising) & $4.99$\\
\bottomrule
\end{tabular}
\end{center}

The framework's distinctive canonical-footing prediction is a \textbf{mild
decline}, $\azero(3)/\azero(0)\approx0.60$--$0.75$, if the DESI-preferred
evolving dark energy is real. If instead $w\to-1$, the framework reduces to
$\Lambda$CDM at all $z$ and the distinctive signal vanishes, leaving only the
$z=0$ tie. The third row is \emph{not} the ALT $z=0$ normalization but a
distinct evolution model in which $\azero$ tracks the total-density Hubble
rate and therefore rises steeply --- the opposite sign. The flat--DESI
separation is only $0.147$ dex at $z=3$ ($0.164$ dex at $z=3.25$); that small
gap is what the galaxy data must resolve.

To be precise about the falsifiable content (and to avoid an overstated
biconditional): $R$ tracks $\sqrt{\rhoDE(z)}$, whatever that history is.
Under the DESI-preferred ($w_0>-1$, $w_a<0$) CPL histories $R$ declines at
high redshift; it reduces to $R=1$ exactly if $w=-1$; and histories with
$\rhoDE$ \emph{rising} toward high $z$ (e.g.\ a constant $w>-1$) would produce
a rising $R$, which under the frozen protocol registers as Verdict~C
(strained) when read at the frozen reference redshift
(Section~\ref{sec:gate}).

\paragraph{The adiabatic-horizon ansatz, stated as an ansatz.} The horizon
relation \eqref{eq:a0} is derived for a \emph{static} de Sitter horizon
(constant $\Lambda$). Promoting it to
$\azero(z)=\tfrac{c}{2}\sqrt{G\rhoDE(z)}$ for an evolving equation of state is
an \textbf{adiabatic / instantaneous-horizon ansatz} --- the assumption that
the inertia scale tracks the instantaneous dark-energy density --- not a
result derived from the de Sitter--Unruh construction, whose horizon is that
of a constant-$\Lambda$ background. The $z$-tracking under test is therefore a
test of this ansatz combined with the measured $w(z)$; a null would disfavour
the ansatz, not the static-horizon relation it reduces to at $z=0$. Relative
to \cite{LPO2008}, who treated $\azero\!-\!H_0$ and $\azero\!-\!\rhoDE$ as two
candidate couplings, the framework's addition is to \emph{fix} the coupling to
$\rhoDE$ via this ansatz (itself posited) --- so the choice between $H_0$ and
$\rhoDE$ is made by the construction, not left free --- and to commit a hashed
decision procedure against the Rubin era.

\paragraph{Non-monotonicity.} Under an evolving CPL history $\azero(z)$ is
\emph{not monotonic}: at the DESI-central rate
($w_0=-0.838$, $w_a=-0.62$) it rises to a small bump --- $R$ peaks at
$\approx1.036$ near $z\approx0.35$ --- and only crosses back below unity at
$z\approx0.9$, declining at higher redshift. (On the DESI+CMB+DESY5 central
history, $w_0=-0.752$, $w_a=-0.86$, the bump is $+6\%$ peaking at
$z\approx0.41$, the phantom-divide crossing \cite{Zdesi}.) The consequences of
summarizing this curve by a single high-$z$ point are drawn out in
Section~\ref{sec:limitations}.

\paragraph{The sharpest galaxy-side signature is a sign.} In the deep-MOND
regime $V^4=GM_{\rm bar}\azero$, so $V\propto \azero^{1/4}$ and
${\rm d}\log V=\tfrac18\,{\rm d}\log\rhoDE$. On the DESY5-central history
($R(3)=0.737$) the framework predicts high-$z$ deep-MOND discs lying $0.033$
dex ($-7.3\%$ in $V$) \emph{below} the local BTFR at $z=3$, where standard
fixed-$\azero$ MOND sits on it and the rival rising model sits $+0.166$ dex
above \cite{Zdesi}. The sign is more robust to amplitude systematics than the
$\azero$ value itself.

\section{The non-circular cross-scale strategy}
\label{sec:noncircular}

$\Lambda$CDM gives no reason for a galaxy's internal acceleration scale to
track cosmic expansion; the framework forces exactly that link. Crucially, the
equality of a galaxy $\azero$ and a cosmic $\Lambda$ at a \emph{single} epoch
is a definition --- one can always convert one into the other, and reading
agreement there is circular. The genuine, non-circular test is whether
\textbf{two independently measured functions of redshift} --- $\azero$ from
galaxy kinematics (parsec--kiloparsec scales) and $\rhoDE$ from cosmic
distances (gigaparsec scales) --- track each other as $z$ varies. The
framework \emph{inherits} the dark-energy history $w(z)$; its content is the
tie between the two datasets, so any decline it exhibits is a consequence of
the framework relation combined with the measured $w(z)$, not an independent
statement about dark energy.

\subsection{Definition versus test, and the $z=0$ anchor}
\label{sec:z0tie}

The $z=0$ equality $\azero(0)=(c/2)\sqrt{G\rho_{\rm DE,0}}$ is definitional: a
$z=0$ match is a consistency knot, not an independent success. Its value is as
the anchor the bridge extends from. The $\Lambda$-blind galaxy $\azero(0)$ is
taken from the gas-dominated SPARC BTFR zero point measured without any
cosmological-$\Lambda$ input (companion measurement \cite{Za0line}):
$\azero(0)=1.181\times10^{-10}\,\msms\,(\pm16\%)$. Against the cosmic-side
$(c/2)\sqrt{G\rho_{\rm DE,0}}$:

\begin{center}
\begin{tabular}{lccc}
\toprule
footing & cosmic $\azero(0)$ & $a_{0,{\rm SPARC}}/a_{0,{\rm cosmic}}$ & tension\\
\midrule
Planck ($\Omega_{\rm DE}=0.685$, $H_0=67.4$) & $9.36\times10^{-11}$ & $1.26$ & $1.44\sigma$\\
SH0ES-$H_0$ ($\Omega_{\rm DE}=0.685$, $H_0=73.0$) & $1.01\times10^{-10}$ & $1.17$ & $0.95\sigma$\\
\bottomrule
\end{tabular}
\end{center}

The tie holds at $\sim1\sigma$ on the canonical (Planck-$H_0$) footing and
slightly better at the local $H_0$. This is by construction and carries no
evidential weight beyond internal consistency; the non-circular content is
entirely in the \emph{slope}.

\subsection{Why a wrong background cannot manufacture the signal}

Because the two datasets never share inputs, a wrong background assumption
shifts both together and cannot manufacture a spurious co-evolution. The mild
distance--cosmology dependence of the galaxy side is quantified:
$\Mbar\propto D^2$, and the ratio $D({\rm DESI})/D({\rm flat})$ shifts $\Mbar$
by $\le0.015$ dex at $z=1$ (smaller at higher $z$) --- roughly $10\times$
below the signal and $15\times$ below the single-object error of
Section~\ref{sec:bigwheel}. Testing whether $\azero$ tracks $\rhoDE$ is not
the same as assuming the DE evolution.

\section{The galaxy-side $a_0(z)$ from the high-$z$ BTFR}
\label{sec:galaxy}

\subsection{The estimator}

On the galaxy side, $\azero(z)$ is read from the BTFR zero point at each
redshift. In the deep-MOND regime $V^4=\azero\,G\Mbar$ is exact (the fitted
slope is forced to $4$), so at fixed circular velocity
\begin{equation}
\frac{\azero(z)}{\azero(0)}=\frac{(V_z/V_0)^4}{M_{{\rm bar},z}/M_{{\rm bar},0}},
\qquad\text{equivalently}\qquad
\Delta\log \azero=-\,\Delta\log \Mbar\big|_{\text{fixed }V}.
\label{eq:btfr}
\end{equation}
Reading $\azero(z)$ is the same operation as reading the BTFR zero point. The
estimator is valid \emph{only} where the fitted slope is genuinely $4$, i.e.\
for objects deep in the low-acceleration regime $g\ll \azero$.

\subsection{The one high-$z$ datum: the Big Wheel at $z=3.25$, now transitional}
\label{sec:bigwheel}

The best available high-$z$ candidate is the ``Big Wheel'' (Wang et al.\ 2025
\cite{Wang2025}; $z=3.2514$), a large low-acceleration disk. A Monte Carlo
over the discovery-paper velocity and baryonic-mass systematics
($V_{\rm rot}=280\pm31\,{\rm km\,s^{-1}}$ with asymmetric-drift correction;
$M_\star=1.7\pm0.7\times10^{11}M_\odot$ dynamically capped;
$M_{\rm gas}=1.8\pm0.9\times10^{11}M_\odot$, $\times1.36$ He) gave
$a_{0,{\rm eff}}=1.54\,(+1.10/-0.61)\times10^{-10}\,\msms$, i.e.\ ratio
$1.31\,(+0.93/-0.52)$ against the SPARC anchor ($1.65$ against canonical,
$1.37$ against ALT --- both footings carried). On those masses the datum was
consistent with a constant and with the DESI decline, and disfavored the
rising $cH(z)$ evolution model ($\sim5\times$) at $\sim2\sigma$.

\textbf{That discovery-mass state is superseded.} A dedicated 2026 ALMA
dynamical model (Quadri et al.\ 2026 \cite{Quadri2026}; W.~Wang among the
authors) finds a mildly rising rotation curve reaching
$V_{\rm rot}\approx314\,{\rm km\,s^{-1}}$ at the outermost \emph{measured}
radius $\approx16.5$ kpc, $M_{\rm gas}=10^{10.76}M_\odot$, and --- critically
--- a stellar mass ambiguous by $0.37$ dex: the dynamical-model fiducial
$M_\star=10^{11.00}M_\odot$ versus a higher SED/half-mass value
$10^{11.37}M_\odot$ that the authors flag as in tension with the dynamics.
Applying the framework's exact kernel
$\azero=(g_{\rm obs}^2-g_{\rm bar}^2)/g_{\rm bar}$ at the measured 16.5-kpc
point, the inferred $\azero$ swings from ratio $0.86$ (SED $M_\star$) to
$3.25$ (dynamical $M_\star$) at central values --- Monte-Carlo medians
$1.28\,[0.41,2.60]$ and $3.11\,[1.89,4.71]$ respectively, with $\sim28\%$ of
the SED-$M_\star$ posterior unphysical ($g_{\rm bar}>g_{\rm obs}$). The object
sits at $y=g_{\rm bar}/\azero\approx0.7$--$1.3$ --- \emph{transitional, not
deep-MOND} --- so $\azero$ is hyper-sensitive to $\Mbar$, and the $0.37$-dex
stellar-mass ambiguity alone spans a factor $\sim3.8$ in $\azero$ ($0.4$--$4.7$
across full systematics). Three consequences, stated plainly:
(i) the single Big Wheel $\azero$ remains $\Mbar$-hostage and cannot separate
the framework's $\sim0.68$ decline at $z=3.25$ from a strong rise;
(ii) the $\sim2\sigma$ exclusion of the rising model held \emph{only on the
discovery masses} --- under the authors' fiducial dynamical mass the object
itself sits at a $\sim3\times$ rise ($3.1\,[1.9,4.7]$), only $\sim1.2\sigma$
from the $\sim5\times$ rising track, so the rise-exclusion too is
$\Mbar$-hostage; and (iii) this is an honest downgrade of the one in-hand
galaxy-side value, and it \emph{empirically confirms} the central finding of
Section~\ref{sec:keyfinding} on the very object one would reach for: with new,
dedicated dynamical data the single-object $\azero$ did not tighten, because
$\Mbar$ calibration, not the object, is decisive. (Scripts:
\texttt{galaxy\_a0z.py} for the discovery-mass Monte Carlo,
\texttt{bigwheel\_update.py} for the v2 re-extraction; Appendix~\ref{app:scripts}.)

\subsection{Why intermediate-$z$ is not a clean $a_0$ probe}
\label{sec:intermediate}

The many $z\sim0.5$--$2.3$ rotating disks (e.g.\ KMOS$^{3\rm D}$) do not help,
for a physical reason: these massive rotators sit at $g\gtrsim \azero$,
\emph{not} deep-MOND, so their fitted BTFR slopes are $3.0$--$3.85$, not $4$,
and the zero point is degenerate with the slope, the pivot, and
baryon/dark-matter-fraction evolution. The inferred ``$\azero$ evolution''
flips sign with the analysis choice: \"Ubler et al.\ (2017)
\cite{Ubler2017} (fixed local slope $3.75$) reads $\azero$ rising
$\sim2$--$3\times$; Sharma et al.\ (2024) \cite{Sharma2024} (free slope
$3.21$) the opposite sign; Di Teodoro et al.\ (2016) \cite{DiTeodoro2016} and
Tiley et al.\ (2019) \cite{Tiley2019} read null. The scatter across analyses
(a factor $\sim2$--$3$, in either direction) dwarfs the $0.70$-vs-$1.0$
signal. These points are not valid $\azero$ readouts; where the $z\approx2.3$
point enters the joint fit below it is placed at its fixed-slope value $1.86$
(rising --- the \emph{wrong} sign for the framework decline), so it counts
against, not for, the framework.

\subsection{The $a_0$--$M_{\rm bar}$ degeneracy}

Because $\Delta\log \azero=-\Delta\log \Mbar$ at fixed $V$, a mis-estimate of
the baryonic mass propagates $1{:}1$ into $\azero$: the two are algebraically
degenerate for this estimator. Rising molecular-gas fractions ($\sim10\%$ at
$z=0$ to $\sim50\%$ at $z=2$), the CO-to-H$_2$ conversion $\alpha_{\rm CO}$,
dust, and IMF evolution all move the BTFR zero point and are indistinguishable
from an $\azero$ shift with the BTFR / point-mass method. This is the crux
limitation developed in Section~\ref{sec:keyfinding}.

\subsection{Does the galaxy data show the decline? No --- neither way}
\label{sec:confront}

Assembling the honest galaxy points (\texttt{confront.py},
Appendix~\ref{app:scripts}):

\begin{center}
\begin{tabular}{lcccc}
\toprule
set & $\chi^2(\text{flat})$ & $\chi^2(\text{DESI-decl})$ & $\Delta\chi^2$ & lean\\
\midrule
all 4 points & $1.08$ & $3.05$ & $-1.97$ & flat, weak\\
$z=0$ + Big Wheel (discovery masses; & $0.27$ & $1.55$ & $-1.28$ & flat, weak\\
\quad transitional per Sec.~\ref{sec:bigwheel}) & & & &\\
\bottomrule
\end{tabular}
\end{center}

The second row is the sole formally meaningful statistic; the ``all 4 points''
row is illustrative only, because the intermediate-$z$ points are declared
invalid $\azero$ readouts (Section~\ref{sec:intermediate}) and cannot
legitimately enter a likelihood. Both rows agree in direction: the data mildly
prefer flat, $|\Delta\chi^2|\approx1.3$--$2.0$ --- not significant --- and
what lean exists in the intermediate-$z$ points is toward \emph{rising}, the
wrong sign for the framework decline. The distinctive $0.60$--$0.75$ decline
is neither detected nor excluded. The separability statistic --- the maximum
sigma at which the current errors could ever tell flat from the DESI decline
--- is $S=0.73\sigma$ (Big Wheel alone), $0.73\sigma$ (clean pair), and
$0.80\sigma$ (all four points): \textbf{underpowered, not passed, not
failed}. Every galaxy point sits $\le1.2\sigma$ from both tracks.

\section{The key finding: $M_{\rm bar}$ calibration, not rotator count, is decisive}
\label{sec:keyfinding}

The decisive question for any future test is how the significance scales. The
naive forecast assumes independent per-object errors that beat down as
$1/\sqrt{N}$: the target signal is $-\log_{10}(0.70)=0.155$ dex, requiring
$\sigma_{\rm mean}=0.052$ dex for $3\sigma$, hence $N\approx4$ rotators at
SPARC-like precision ($0.10$ dex) or $N\approx24$ at realistic high-$z$
per-object precision ($0.25$ dex). \textbf{This is wrong}, because the
dominant $\Mbar$ error --- $\alpha_{\rm CO}$, IMF, and dust prescriptions
applied \emph{identically} to a sample --- is largely common-mode, and a
common-mode error does not average down with $N$. With
$\sigma_{\rm mean}^2=\sigma_{\rm ind}^2/N+\sigma_{\rm cm}^2$ and a correlated
$\Mbar$ floor $\sigma_{\rm cm}$:

\begin{center}
\begin{tabular}{lc}
\toprule
$\Mbar$ common-mode floor & $N$ for $3\sigma$\\
\midrule
$0.00$ dex & $24$\\
$0.05$ dex & $\sim377$\\
$\ge0.08$ dex & \textbf{never reaches $3\sigma$ at any $N$}\\
\bottomrule
\end{tabular}
\end{center}

Because realistic high-$z$ $\Mbar$ calibration floors are $\sim0.08$--$0.15$
dex ($\alpha_{\rm CO}$ alone is a $0.1$--$0.2$ dex systematic), the test is
decisive \textbf{only if the common-mode $\Mbar$ systematic can be driven
below $\sim0.05$ dex} --- a calibration problem, not a counting problem. The
2026 ALMA re-analysis of the Big Wheel (Section~\ref{sec:bigwheel})
demonstrates this on real data: a $0.37$-dex stellar-mass ambiguity on a
single dedicated-data object swings its $\azero$ by a factor $\sim3.8$ ---
exactly the $\Mbar$-hostage behaviour this section predicts, and precisely why
more objects at fixed $\Mbar$ calibration would not help. This is the paper's
core methodological contribution: \textbf{the decisive observational lever is
a uniformly calibrated $\Mbar$ ladder to $\lesssim0.05$ dex common-mode, not
the number of rotators.} (Both footings enter identically --- the forecast is
in ratios and dex, so the $20.9\%$ footing choice cancels.)

\subsection{Forecast requirements, correctly stated}
\label{sec:forecastreq}

Reaching $3\sigma$ on the $0.60$--$0.75$ decline versus flat therefore
requires \emph{both} (i) a common-mode $\Mbar$ calibration $\lesssim0.05$ dex
\emph{and} (ii) a sample of $\sim20$--$40$ clean deep-MOND $z\sim2$--$3$
rotators, Big-Wheel-like, with gas-traced outer rotation curves (JWST/ALMA) or
individually resolved deep-MOND curves (ELT/HARMONI); the per-object precision
required is $\le0.16$ dex ($N=10$), $\le0.23$ dex ($N=20$), $\le0.33$ dex
($N=40$), all under the independence assumption that returns the argument to
the common-mode floor above. Neither condition alone suffices: without the
calibration floor, no rotator count reaches $3\sigma$. Clean deep-MOND disks
like the Big Wheel are rare, so assembling $20$--$40$ is itself an
observational stretch.

\subsection{The $a_0$-separable observable}

The deeper methodological escape is to stop using the BTFR zero point at all.
Because $\azero$ and $\Mbar$ are $1{:}1$ degenerate in the BTFR, no amount of
BTFR data isolates $\azero$ from baryonic-mass evolution. The
\textbf{$\azero$-separable} observable is the \emph{shape} of the rotation
curve --- the radial-acceleration-relation (RAR) transition --- which depends
on $\azero$ at fixed $\Mbar$ and therefore breaks the degeneracy the zero
point cannot. A high-$z$ RAR-transition measurement reads $\azero$ directly,
without hostage to the gas-mass normalization. This is the observable to
target: not more BTFR zero points, but resolved high-$z$ rotation-curve
shapes.

\section{The cosmology-side gate and the pre-registration}
\label{sec:gate}

\subsection{DESI DR2 propagated through the relation}
\label{sec:desi}

On the cosmic side there is no free dark-energy parameter to fit --- the
framework consumes whatever $\rhoDE(z)$ the data prefer. Propagating the DESI
DR2 \cite{DESI2025} $w_0w_a$CDM posterior (BAO + CMB + SNe; representative
marginals with $w_0$--$w_a$ correlation $\approx-0.87$; $\Lambda$CDM excluded
at $2.8\sigma$/Pantheon+, $4.2\sigma$/DESY5, $3.8\sigma$/Union3) through
Eq.~\eqref{eq:law} (\texttt{desi\_posterior\_a0z.py},
Appendix~\ref{app:scripts}):

\begin{center}
\begin{tabular}{lccccc}
\toprule
SNe combination & $\azero(3)/\azero(0)$ & $[16\%,84\%]$ & $P({\rm declines})$ & signif. & in $0.60$--$0.75$\\
\midrule
DESI+CMB+Pantheon+ & $0.775$ & $[0.683,\,0.878]$ & $97.9\%$ & $2.0\sigma$ & $38\%$\\
DESI+CMB+DESY5 & $0.737$ & $[0.651,\,0.833]$ & $99.3\%$ & $2.5\sigma$ & $51\%$\\
DESI+CMB+Union3 & $0.707$ & $[0.599,\,0.831]$ & $98.3\%$ & $2.1\sigma$ & $48\%$\\
\bottomrule
\end{tabular}
\end{center}

Across the three combinations the median is $\azero(3)/\azero(0)\simeq0.74$
with a $\sim2.2\sigma$ decline, and all three posteriors overlap the
framework's $0.60$--$0.75$ band. This is a genuine, live cosmology-side hint:
\emph{if} the DESI evolving-dark-energy signal is real, the induced
$\azero(z)$ decline falls inside the framework band. It is not, however, an
independent test --- it is a re-expression of the DESI $w(z)$ result through
the framework relation, the same DESI information seen twice, and it cannot be
added to the galaxy-side $\chi^2$ as if it were a second measurement. It
dissolves entirely if the DESI signal regresses to $w=-1$.

\subsection{The frozen estimator and pre-committed thresholds}
\label{sec:frozen}

The decision procedure was frozen, SHA-256-hashed, and Zenodo-timestamped on
2026-07-21 \cite{Zprereg}, before any calibrated Rubin supernova cosmology
exists (LSST began operations 2026-06-30; the LSST-DESC calibrated Type-Ia
sample is not expected before $\sim$2027). The estimator
(\texttt{a0z\_gate\_estimator.py}) draws $N=4\times10^5$ samples from the
input $(w_0,w_a)$ bivariate Gaussian, maps each through Eq.~\eqref{eq:law} to
$R=\azero(3)/\azero(0)$, and reports the median, $[16,84]$ interval, and the
posterior credibility of a decline mapped to a Gaussian sigma. \textbf{No
threshold, band, or estimator may change after the freeze; only the input
$(w_0,w_a)$ is updated.}

\begin{center}
\begin{tabular}{p{4.2cm}p{9.3cm}}
\toprule
Verdict & Pre-committed condition\\
\midrule
\textbf{A) GATE OPEN} --- distinctive prediction alive and supported &
decline at $\ge3\sigma$ \textbf{and} $R\in[0.55,0.85]$ $\to$ $\azero$ evolves
as required; warrants the independent galaxy-side confrontation. GATE OPEN
means ``survives to face the galaxy-side test,'' explicitly not
``confirmed.''\\
\textbf{B) GATE DISSOLVED} --- safe core, \emph{not} falsification &
consistent with flat ($R=1$ within $2\sigma$) $\to$ framework reduces to
constant-$\azero$ MOND; the distinctive prediction is inapplicable.\\
\textbf{C) FRAMEWORK STRAINED} --- cosmology-side tension &
$R>1$ (a rise) at $\ge3\sigma$ $\to$ the scale rises with $z$, which
$\azero\propto\sqrt{\rhoDE}$ does not permit under evolving DE.\\
\textbf{UNDECIDED} & none of the above met.\\
\bottomrule
\end{tabular}
\end{center}

The band $[0.55,0.85]$ is fixed and hashed. Its width comes from the spread of
currently-allowed dark-energy histories --- holding $w_a=-0.62$:
$w_0\approx-0.79\to R=0.85$ (mild evolution), $w_0\approx-1.00\to R=0.55$
(strong evolution) --- \emph{not} from the two value footings (which give an
identical $R$ at fixed $(w_0,w_a)$ and cannot widen it). Two honesty notes:
(i) the band is deliberately wider than the nominal distinctive claim
$R\approx0.60$--$0.75$, and the DESI value $0.775$ sits just above the nominal
$0.75$ though inside the operational band; (ii) the upper edge $0.85$ is close
to flat, so ``supported'' at that edge is a weak decline --- the intended cost
of a band broad enough not to be tuned to one posterior. A decline steeper
than $R<0.55$ at $\ge3\sigma$ returns UNDECIDED, not a verdict --- the frozen
procedure has no ``over-decline strain'' branch mirroring Verdict~C.

\subsection{Data hierarchy and the refused DIY-Hubble-diagram path}

The gate is fed by a strict, frozen hierarchy: (1) \textbf{primary} --- the
LSST-DESC \emph{calibrated} Type-Ia sample ($\sim$2027+), the only input that
can produce a supported verdict; (2) \textbf{interim} --- published Rubin
$w_0w_a$ constraints as they appear; (3) \textbf{readiness tracker only} ---
a running count of broker-classified Type-Ia candidates from the public alert
stream (\texttt{broker\_sn\_counter.py}), which \emph{never} feeds the
verdict. We explicitly refuse to build a do-it-yourself Hubble diagram from
public broker light curves: without the DESC calibration layer (zero-point
stability, photometric-redshift systematics, modeled selection function), any
$w(z)$ extracted from raw broker photometry is systematics-dominated, and a
``Rubin supports $\azero(z)$'' claim built on it would be a manufactured
result.

\subsection{State at the freeze (2026-07-21)}

Running the frozen estimator on the best current input --- DESI DR2
\cite{DESI2025}, $w_0=-0.838\pm0.055$, $w_a=-0.62\pm0.22$, correlation
$-0.86$ (the DESI+CMB+Pantheon+ combination, the mildest; DESY5 and Union3
prefer stronger evolution and the verdict is unchanged under them) --- gives
$R=0.775\,[0.683,0.878]$, a decline at $2.0\sigma$: below the committed
$3\sigma$ threshold, hence \textbf{UNDECIDED} --- a hint, explicitly not a
detection, explicitly not gate-open. The broker readiness tracker stood at
36,836 Type-Ia-classified objects at the freeze; by the frozen rule it feeds
timing only, never the verdict.

\section{Forecast against Rubin/LSST}
\label{sec:forecast}

The committed forecast (\texttt{forecast\_rubin\_a0z.py}) propagates
representative projected Rubin $w_0$--$w_a$ covariances (photometric SN
figure-of-merit $\sim$150 alone, $\sim$500 with CMB/BAO \cite{DESC,Ivezic2019})
through Eq.~\eqref{eq:law}, taking the DESI-central evolving history as the
assumed truth:

\begin{center}
\begin{tabular}{lcccc}
\toprule
Input scenario & $\sigma(w_0)$ & $\sigma(w_a)$ & $R=\azero(3)/\azero(0)$ & Decline signif.\\
\midrule
DESI DR2 now (baseline) & 0.055 & 0.22 & 0.775 [0.683, 0.878] & $2.0\sigma$\\
Rubin SN alone (FoM $\sim$150) & 0.040 & 0.15 & 0.775 [0.717, 0.837] & $3.3\sigma$\\
Rubin SN + CMB/BAO (FoM $\sim$500) & 0.020 & 0.08 & 0.775 [0.742, 0.809] & $4.6$--$4.7\sigma$\\
\bottomrule
\end{tabular}
\end{center}

The \textbf{mirror test}: if the universe is in fact flat ($w=-1$, $\azero$
truly constant), Rubin-alone pins $R=1.00\pm0.078$, excluding $0.74$ at
$3.3\sigma$; Rubin + CMB/BAO gives $R=1.00\pm0.043$, excluding it at
$6.0\sigma$. The honest summary is deliberately unheroic: if dark energy
evolves at the DESI-central rate, Rubin sharpens the decline from
$\sim2.0\sigma$ to $\sim3.3\sigma$ (SNe alone) or $\sim4.6$--$4.7\sigma$ (with
CMB/BAO); if dark energy is flat, Rubin excludes the framework's $0.74$ at
$\sim3.3$--$6.0\sigma$. Either way the gate is likely settled at the
$3$--$6\sigma$ level by roughly 2028--2030 --- a real upgrade on DESI-now, not
a one-instrument slam dunk. These are projections, not data, and cannot
produce a real verdict.

\section{Limitations}
\label{sec:limitations}

Because the estimator and thresholds are hashed and cannot be revised after
the freeze, the following limitations are documented rather than fixed. They
are stated here so that no reading can present the test as tighter than it is.

\begin{enumerate}
\item \textbf{The signal is a $z=3$ extrapolation beyond the supernova range,
and the significance is $z$-dependent.} The LSST photometric SN sample
constrains $w(z)$ mainly at $z\lesssim1.2$; the gate is read at $z=3$.
Propagating the \emph{same} DESI-DR2 posterior to different reference
redshifts gives a decline significance of $0.34\sigma$ at $z=1$ ($R=0.99$,
effectively flat), $1.7\sigma$ at $z=2$, $2.0\sigma$ at $z=3$, and $2.4\sigma$
at $z=10$; for the tight SN+CMB/BAO forecast covariance it is $1.1\sigma$ at
$z=1$ but $4.7\sigma$ at $z\ge2$ (reproduced by
\texttt{significance\_vs\_z.py}, Appendix~\ref{app:scripts}). What the gate
genuinely tests is whether the SN-constrained posterior excludes
$(w_0,w_a)=(-1,0)$; the $z=3$ evaluation converts that into an $\azero$-ratio
but adds no independent $z=3$ information. $R(z{=}3)=0.775$ must \emph{not} be
read as a measured $z=3$ acceleration scale.
\item \textbf{The framework's own $\azero(z)$ bumps up at low redshift, which
is where the galaxy-side test lives.} At the DESI-central history $R$ peaks at
$\approx1.036$ near $z\approx0.35$ and only falls below $1$ beyond
$z\approx0.9$. A GATE-OPEN verdict on the $z=3$ decline does not translate
into a low-$z$ decline for the galaxy side to corroborate; at $z\lesssim0.5$
the framework predicts a slight \emph{rise} (a few percent), at $z\sim1$
essentially no change. Anyone carrying ``the framework predicts $\azero$
declines'' to low-$z$ observers would have the sign wrong. The two lanes probe
opposite branches of the same non-monotonic curve.
\item \textbf{Verdict C is well-posed only at the frozen $z=3$.} Because
$\azero(z)$ genuinely rises for $z\lesssim0.9$, a measured rise contradicts
the $\sqrt{\rhoDE}$ footing only when read at high $z$; Verdict~C must never
be applied to a low-$z$ input.
\item \textbf{Doc-vs-code items in the frozen protocol.} (a) The Verdict-B
sub-clause (``if $R\le0.85$ is additionally excluded at $\ge3\sigma$, the
decline is positively ruled out'') is narrative only; the frozen script
computes Verdict~B purely as consistency with $R=1$ within $2\sigma$. (b) The
over-decline gap ($R<0.55$ significant $\to$ UNDECIDED) is a real asymmetry
with Verdict~C. (c) The frozen estimator prints GATE OPEN for the two
\emph{forecast} covariances; those are fabricated inputs shown for
illustration and, by the data hierarchy, can never constitute a verdict.
\item \textbf{Significance convention.} The reported ``$N\sigma$'' is a
Bayesian posterior tail (the fraction of samples with $R\ge1$ mapped to a
Gaussian sigma), not a frequentist likelihood-ratio detection; the estimator
($N=4\times10^5$) and forecast script ($N=3\times10^5$) differ at the
$0.1\sigma$ level ($4.7$ vs $4.6\sigma$), reported as the range.
\item \textbf{$\Mbar$ dominance of the galaxy side.} The independent lane is
currently underpowered ($S\simeq0.8\sigma$), $\Mbar$-hostage
(Sections~\ref{sec:bigwheel}--\ref{sec:keyfinding}), and its intermediate-$z$
literature is sign-contested.
\item \textbf{The covariant completion is not in hand, and the test does not
lean on it.} The observable $R$ requires only the scalar relation
Eq.~\eqref{eq:law} and a $w(z)$ posterior --- it does not invoke a Lagrangian.
A separate modified-inertia field-theory program exists \cite{Zmift} but this
test is scoped so that nothing rests on it; the absence of a completed
covariant modified-inertia theory is a limitation of the framework, disclosed,
not of the test.
\item \textbf{The adiabatic-horizon ansatz is posited}
(Section~\ref{sec:tracks}), and the evolution-law fork
(Section~\ref{sec:forks}) is resolved by pre-commitment, not by derivation.
\end{enumerate}

\section{Conclusion}

We have set out, and pre-registered, a non-circular cross-scale test of de
Sitter--Unruh modified inertia: the framework replaces $\Lambda$CDM's free
dark-energy fit with the galaxy acceleration scale, forcing
$\azero(z)/\azero(0)=\sqrt{\rhoDE(z)/\rho_{\rm DE,0}}$ --- a bridge between
galaxy kinematics and cosmic distances that $\Lambda$CDM structurally lacks.
The $z=0$ tie holds at $\sim1\sigma$ (the definitional anchor). The cosmology
side is live but inherited: the DESI DR2 posterior propagated through the
relation yields $\azero(3)/\azero(0)\simeq0.74$ at $\sim2.2\sigma$ --- below
the committed $3\sigma$ bar, hence UNDECIDED --- and it dissolves if
$w\to-1$. The galaxy side is underpowered: the distinctive $0.60$--$0.75$
decline is neither detected nor excluded ($S\simeq0.8\sigma$); the one
$z=3.25$ object is transitional and $\Mbar$-hostage across ratio $0.4$--$4.7$,
and even its former $\sim2\sigma$ exclusion of the rising model held only on
the discovery masses. The central finding is methodological and redirects the
observational strategy: the dominant $\Mbar$ systematic is common-mode, so
with a $\gtrsim0.08$-dex correlated floor the BTFR test never reaches
$3\sigma$ at any rotator count; the decisive lever is $\Mbar$ calibration to
$\lesssim0.05$ dex common-mode, and the $\azero$-separable observable is the
rotation-curve/RAR transition, not the BTFR zero point. The estimator,
thresholds, data hierarchy, and refusal of the DIY-Hubble path are hashed and
timestamped before the deciding Rubin data exist; forecasts project a
$3$--$6\sigma$ settlement by $\sim$2028--2030. The value of $\azero$, the
normalization $Z$, and the sign of the effect are posited, not derived; both
footings are carried throughout; and the contribution is a bounded,
falsifiable, pre-registered program, not a detection.

\section*{Data and code availability}

Every load-bearing quantitative claim in this manuscript traces to a
committed, runnable, exit-0 Python script (numpy/scipy) in the author's
public repository, which is \emph{self-verifying}: re-running the committed
scripts reproduces every number in the text and tables
(Appendix~\ref{app:scripts} lists the inventory and the freeze hashes). The
verification toolkit is released as \texttt{a0kit} (Zenodo software concept
DOI \href{https://doi.org/10.5281/zenodo.21478981}{10.5281/zenodo.21478981};
an ASCL submission is pending editor review). The consolidated source deposits are Zenodo DOIs
10.5281/zenodo.20737162 (the $\azero(z)$ law and DESI confrontation
\cite{Zdesi}), 10.5281/zenodo.21440407 (the cross-scale test \cite{Zcross}),
and 10.5281/zenodo.21478568 (the Rubin/LSST pre-registration \cite{Zprereg});
companions are 10.5281/zenodo.21403470 (field theory \cite{Zmift}) and
10.5281/zenodo.21419735 (the SPARC $\azero$-line measurement \cite{Za0line}).

\appendix

\section{Freeze integrity and script inventory}
\label{app:scripts}

\subsection*{Freeze hashes}

The SHA-256 hashes of the frozen artifacts, recorded in
\texttt{FREEZE\_HASHES.txt} and verified 2026-07-21 \cite{Zprereg}:

{\small
\begin{verbatim}
852ca954431b9239415038dde9f2b2ad9f89fe6325ff7a8ba3e3ebe357366740  a0z_gate_estimator.py
ba24927ddb3263c390646019b1dba1968c4aa3de5ed572bc5de57685afd2b77f  broker_sn_counter.py
bd714c93edd3eadc6c84413dee331e86be7ebaca8f0a5f88c84ec1542f46df0f  PREREGISTRATION.md
530e66beef425bb500d085051c3219669e6d6e31364eec9263ea036bca7a5697  forecast_rubin_a0z.py
\end{verbatim}
}

\noindent No threshold, band, or estimator changes after these hashes; only
the calibrated $(w_0,w_a)$ input is updated when the LSST-DESC sample lands.
The Zenodo deposit timestamp \cite{Zprereg} is the public lock-in.

\subsection*{Script inventory}

All numbers in this manuscript are reproduced by the following committed
scripts (numpy/scipy only, exit 0, both footings carried):

\begin{itemize}
\item \texttt{a0z\_gate\_estimator.py} (frozen) --- the gate: verdict logic,
thresholds, $R$ posterior; reproduces $R=0.775\,[0.683,0.878]$, $2.0\sigma$,
UNDECIDED.
\item \texttt{forecast\_rubin\_a0z.py} (frozen) --- the Rubin forecast table
and the mirror test.
\item \texttt{broker\_sn\_counter.py} (frozen) --- the readiness tracker
(timing only; never feeds the verdict).
\item \texttt{desi\_posterior\_a0z.py} --- the three-SNe-combination DESI
propagation table of Section~\ref{sec:desi}.
\item \texttt{galaxy\_a0z.py} --- the galaxy-side ladder: SPARC anchor,
sign-contested intermediate-$z$ compilation, discovery-mass Big-Wheel Monte
Carlo.
\item \texttt{bigwheel\_update.py} --- the v2 ALMA-dynamical-model
re-extraction of the Big Wheel $\azero$ (both stellar-mass scenarios; the
$0.4$--$4.7$ ratio swing of Section~\ref{sec:bigwheel}).
\item \texttt{confront.py} --- the joint $\chi^2$, separability $S$, $z=0$
tie, and the common-mode-floor forecast table of
Section~\ref{sec:keyfinding}.
\item \texttt{significance\_vs\_z.py} (post-freeze traceability addendum;
deliberately \emph{not} in the freeze hashes) --- the $z$-dependent decline
significances of Section~\ref{sec:limitations}, item 1.
\end{itemize}

Supporting reconstructions from an independent lane confirm that Type Ia
supernovae alone cannot measure $\azero(z{=}3)$: the dark-energy residual
$E^2-\Omega_m(1+z)^3$ is a difference of large numbers, and the
model-independent Gaussian-process reconstruction \cite{Seikel2012} is
sign-pinned only to $z\approx0.40$ --- which is exactly why the two-dataset
(galaxy $\times$ cosmic-distance) construction, and not SNe alone, is the
test.

\begin{thebibliography}{99}

\bibitem{Milgrom1983} M.~Milgrom,
\emph{A modification of the Newtonian dynamics as a possible alternative to
the hidden mass hypothesis}, Astrophys.\ J.\ \textbf{270} (1983) 365.

\bibitem{Milgrom1999} M.~Milgrom,
\emph{The modified dynamics as a vacuum effect},
Phys.\ Lett.\ A \textbf{253} (1999) 273 [astro-ph/9805346] ---
the $\nu=\sqrt{1+1/y}$ kernel (Eq.~9).

\bibitem{LPO2008} M.~Limbach, D.~Psaltis and F.~\"Ozel,
\emph{The Redshift Evolution of the Tully--Fisher Relation as a Test of
Modified Gravity}, arXiv:0809.2790 (2008) --- the $\azero\sim c\sqrt{\rhoDE}$
coincidence and the $\azero$--$H_0$/$\azero$--$\rhoDE$ couplings tested
against Tully--Fisher data to $z\simeq1.2$.

\bibitem{AeST} C.~Skordis and T.~Zlosnik,
\emph{New Relativistic Theory for Modified Newtonian Dynamics},
Phys.\ Rev.\ Lett.\ \textbf{127} (2021) 161302.

\bibitem{DESI2025} DESI Collaboration,
\emph{DESI DR2 Results II: Measurements of Baryon Acoustic Oscillations and
Cosmological Constraints}, arXiv:2503.14738; Phys.\ Rev.\ D \textbf{112}
(2025) 083515.

\bibitem{Brout2022} D.~Brout et al.,
\emph{The Pantheon+ Analysis: Cosmological Constraints},
Astrophys.\ J.\ \textbf{938} (2022) 110.

\bibitem{Lelli2016} F.~Lelli, S.~S.~McGaugh and J.~M.~Schombert,
\emph{SPARC: Mass Models for 175 Disk Galaxies with Spitzer Photometry and
Accurate Rotation Curves}, Astron.\ J.\ \textbf{152} (2016) 157.

\bibitem{McGaugh2005} S.~S.~McGaugh,
\emph{The Baryonic Tully--Fisher Relation of Galaxies with Extended Rotation
Curves and the Stellar Mass of Rotating Galaxies},
Astrophys.\ J.\ \textbf{632} (2005) 859.

\bibitem{McGaugh2012} S.~S.~McGaugh,
\emph{The Baryonic Tully--Fisher Relation of Gas-rich Galaxies as a Test of
$\Lambda$CDM and MOND}, Astron.\ J.\ \textbf{143} (2012) 40.

\bibitem{DiTeodoro2016} E.~M.~Di Teodoro, F.~Fraternali and S.~H.~Miller,
\emph{Flat rotation curves and low velocity dispersions in KMOS star-forming
galaxies at $z\sim1$}, Astron.\ Astrophys.\ \textbf{594} (2016) A77.

\bibitem{Ubler2017} H.~\"Ubler et al.,
\emph{The Evolution of the Tully--Fisher Relation between $z\sim2.3$ and
$z\sim0.9$ with KMOS$^{3\rm D}$}, Astrophys.\ J.\ \textbf{842} (2017) 121.

\bibitem{Tiley2019} A.~L.~Tiley et al.,
\emph{The shapes of the rotation curves of star-forming galaxies over the
last $\approx$10 Gyr}, Mon.\ Not.\ Roy.\ Astron.\ Soc.\ \textbf{482} (2019)
2166.

\bibitem{Sharma2024} G.~Sharma, V.~Upadhyaya, P.~Salucci and S.~Desai,
\emph{Tully--Fisher relation at $0.6\le z\le2.5$},
Astron.\ Astrophys.\ \textbf{690} (2024) [arXiv:2406.08934] --- free bTFR
slope $3.21\pm0.28$.

\bibitem{Wang2025} W.~Wang et al.,
\emph{A giant disk galaxy two billion years after the Big Bang (the ``Big
Wheel'')}, Nature Astron.\ \textbf{9} (2025) 710 [arXiv:2409.17956].

\bibitem{Quadri2026} G.~Quadri, S.~Cantalupo, C.~Bacchini et al.,
\emph{The galaxy--halo connection and dynamical model of the Big Wheel giant
disc at $z\simeq3.25$}, Astron.\ Astrophys.\ (2026), in press
[arXiv:2605.04144].

\bibitem{Seikel2012} M.~Seikel, C.~Clarkson and M.~Smith,
\emph{Reconstruction of dark energy and expansion dynamics using Gaussian
processes}, JCAP \textbf{06} (2012) 036.

\bibitem{Ivezic2019} \v{Z}.~Ivezi\'c et al.\ (LSST/Rubin),
\emph{LSST: From Science Drivers to Reference Design and Anticipated Data
Products}, Astrophys.\ J.\ \textbf{873} (2019) 111.

\bibitem{DESC} LSST Dark Energy Science Collaboration,
supernova cosmology forecasts (photometric FoM $\sim$150; with CMB/BAO
$\sim$500).

\bibitem{Zdesi} C.~P.~Zimmerman,
\emph{Dark Energy Sets the Galaxy Acceleration Scale: A Parameter-Free
$\azero(z)$ Test $\Lambda$CDM Cannot Take}, Zenodo (2026),
DOI \href{https://doi.org/10.5281/zenodo.20737162}{10.5281/zenodo.20737162}
(author's preprint deposit; consolidated here).

\bibitem{Zcross} C.~P.~Zimmerman,
\emph{Does the Galaxy Acceleration Scale Track Cosmic Dark Energy? A
Non-Circular Cross-Scale $\azero(z)$ Test}, Zenodo (2026), concept DOI
\href{https://doi.org/10.5281/zenodo.21440407}{10.5281/zenodo.21440407}
(author's preprint deposit, v2; consolidated here).

\bibitem{Zprereg} C.~P.~Zimmerman,
\emph{A Pre-Registered $\azero(z)$ Gate for the Rubin/LSST Supernova Stream:
Committing the Test Before the Data}, Zenodo (2026),
DOI \href{https://doi.org/10.5281/zenodo.21478568}{10.5281/zenodo.21478568}
(author's pre-registration deposit; consolidated here).

\bibitem{Zmift} C.~P.~Zimmerman,
\emph{MI Field Theory Results 2026}, Zenodo (2026),
DOI \href{https://doi.org/10.5281/zenodo.21403470}{10.5281/zenodo.21403470}
(companion; cited, not consolidated).

\bibitem{Za0line} C.~P.~Zimmerman,
\emph{Reading the Cosmological Constant from Dwarf-Galaxy Rotation Curves:
The $\azero$-Line}, Zenodo (2026),
DOI \href{https://doi.org/10.5281/zenodo.21419735}{10.5281/zenodo.21419735}
(companion; cited, not consolidated).

\end{thebibliography}

\end{document}
